\documentclass[aspectratio=169,10pt,t]{beamer}
\usepackage{ac1}
\aulainfo{06}{Estrutura Básica de Programas}{Sequência, entrada e saída, expressões simples}

\begin{document}
\capa

\begin{frame}{Roteiro}
  \begin{itemize}
    \item Programação estruturada: três blocos de construção (Cap.~3.1)
    \item \textbf{Sequência} (3.1.1): executar comandos em ordem
    \item Saída com \texttt{print}: separadores e formatação
    \item Entrada com \texttt{input} e conversão de tipos
    \item Expressões e precedência de operadores
    \item \textbf{Aula prática} no laboratório
  \end{itemize}
\end{frame}

\begin{frame}{Programação estruturada}
  Todo programa pode ser construído com apenas três \textbf{blocos}:
  \vspace{1ex}
  \begin{columns}[T]
    \column{0.32\textwidth}
    \begin{block}{Sequência}
      Um comando depois do outro, de cima para baixo.
    \end{block}
    \column{0.32\textwidth}
    \begin{block}{Seleção}
      Decidir entre caminhos: \texttt{if}/\texttt{else} (Aula 7).
    \end{block}
    \column{0.32\textwidth}
    \begin{block}{Iteração}
      Repetir um bloco: \texttt{while}/\texttt{for} (Aula 8).
    \end{block}
  \end{columns}
  \vspace{2ex}
  Hoje: \textbf{sequência}. A ordem importa: o computador nunca ``adivinha'' o que vem depois.
\end{frame}

\begin{frame}[fragile]{Sequência: a ordem dos comandos}
\begin{lstlisting}
base = 5
altura = 6
area_triang = (base * altura) / 2
print('A área do triângulo com base', base,
      'e altura', altura, 'é', area_triang)
\end{lstlisting}
\begin{lstlisting}[style=saida]
A área do triângulo com base 5 e altura 6 é 15.0
\end{lstlisting}
  \begin{alertblock}{O que acontece se invertermos as linhas 3 e 4?}
    \texttt{NameError}: a variável \texttt{area\_triang} ainda não existe quando \texttt{print} é executado.
  \end{alertblock}
  \fonte{code/Cap2/prog.2.5\_area\_triangulo\_v3.py}
\end{frame}

\begin{frame}[fragile]{Evolução do programa: três versões}
\begin{lstlisting}
# v1: um valor "mágico" (sem contexto)
print((5*6)/2)

# v2: guarda o resultado e explica a saída
area_triang = (5*6)/2
print('A área do triângulo é ', area_triang)

# v3: dados em variáveis com nomes claros
base = 5
altura = 6
\end{lstlisting}
  \begin{itemize}
    \item Cada versão é mais fácil de \alert{ler e modificar} que a anterior
    \item Evitar números ``soltos'' no meio do código
  \end{itemize}
  \fonte{code/Cap2/prog.2.3, prog.2.4, prog.2.5}
\end{frame}

\begin{frame}[fragile]{Saída: a função \texttt{print}}
\begin{lstlisting}
print("Olá", "mundo")             # Olá mundo
print("Olá", "mundo", sep="-")    # Olá-mundo
print("a", end=" ")               # não pula linha
print("b")                        # a b
print("Total:", 7 * 3)            # Total: 21
\end{lstlisting}
  \begin{itemize}
    \item Vários argumentos separados por vírgula; por padrão, separados por um espaço
    \item \texttt{sep} troca o separador; \texttt{end} troca o final (padrão: nova linha)
    \item \texttt{\textbackslash n} = nova linha, \texttt{\textbackslash t} = tabulação
  \end{itemize}
  \fonte{code/Cap3/prog.3.15\_contagem.py (usa \texttt{end=' '})}
\end{frame}

\begin{frame}[fragile]{Saída formatada}
\begin{lstlisting}
import math
x = math.pi
print("Pi vale", x)
print("Pi vale {:.2f}".format(x))      # duas casas decimais
print(f"Pi vale {x:.4f}")               # f-string
print(f"{'Graus':>9}{'Seno':>9}")       # alinha à direita em 9 colunas
print(f"{30:>9.2f}{math.sin(math.radians(30)):>9.2f}")
\end{lstlisting}
\begin{lstlisting}[style=saida]
Pi vale 3.141592653589793
Pi vale 3.14
Pi vale 3.1416
    Graus     Seno
    30.00     0.50
\end{lstlisting}
  \fonte{code/Cap6/prog.6.17\_tabela\_senos\_arquivo\_f.py}
\end{frame}

\begin{frame}[fragile]{Entrada: a função \texttt{input}}
\begin{lstlisting}
nome = input("Qual é o seu nome? ")
print("Olá,", nome)
idade = input("Qual a sua idade? ")
print(type(idade))             # <class 'str'>
\end{lstlisting}
  \begin{alertblock}{\texttt{input} sempre devolve \texttt{str}!}
    Para fazer contas é preciso converter: \texttt{int(\ldots)} ou \texttt{float(\ldots)}.
  \end{alertblock}
  \begin{itemize}
    \item O texto entre parênteses é o \emph{prompt} exibido ao usuário
    \item Se o usuário digitar algo inválido (\texttt{abc} para \texttt{int}), ocorre \texttt{ValueError}; trataremos na Aula 13
  \end{itemize}
\end{frame}

\begin{frame}[fragile]{O erro clássico: somar textos}
\begin{lstlisting}
nota1 = input("Qual é a primeira nota? ")
nota2 = input("E a segunda nota? ")
nota_total = nota1 + nota2
print("A soma é ", nota_total)
\end{lstlisting}
\begin{lstlisting}[style=saida]
Qual é a primeira nota? 7
E a segunda nota? 8
A soma é  78
\end{lstlisting}
  \begin{itemize}
    \item \texttt{'7' + '8'} \(\Rightarrow\) \texttt{'78'} (concatenação de strings)
    \item \textbf{Correção:} \texttt{nota1 = float(input(\ldots))}
  \end{itemize}
  \fonte{code/Cap3/prog.3.2\_testa\_input.py}
\end{frame}

\begin{frame}[fragile]{Programa completo: área de um quadrado}
\begin{lstlisting}
# lê o valor do lado
lado = float(input("Digite o valor do lado do quadrado: "))
# calcula a área
area = lado * lado
# imprime o resultado
print("Área do quadrado:", area)
\end{lstlisting}
  Estrutura típica de um programa simples: \textbf{entrada \(\to\) processamento \(\to\) saída}.
  \begin{exampleblock}{Exercícios 3.1 e 3.2}
    Versão com \texttt{int} e versão com \texttt{float}: qual a diferença para o lado \texttt{2.5}?
  \end{exampleblock}
\end{frame}

\begin{frame}{Expressões e precedência}
  \centering\small
  \begin{tabular}{cl}
    \toprule
    Ordem & Operadores\\
    \midrule
    1 & parênteses \texttt{( )}\\
    2 & potência \texttt{**}\\
    3 & \texttt{+x}, \texttt{-x} (sinal)\\
    4 & \texttt{*}, \texttt{/}, \texttt{//}, \texttt{\%}\\
    5 & \texttt{+}, \texttt{-}\\
    6 & comparações \texttt{==}, \texttt{<}, \ldots\\
    7 & \texttt{not}, \texttt{and}, \texttt{or}\\
    \bottomrule
  \end{tabular}

  \vspace{1.5ex}\normalsize
  \(2 + 3 \times 4 = 14\), mas \((2+3)\times4 = 20\). Na dúvida, \alert{use parênteses}.
  \vspace{1ex}

  \footnotesize Cuidado: \(\dfrac{a+b}{2}\) é \texttt{(a + b) / 2}, e não \texttt{a + b / 2}.
\end{frame}

\begin{frame}[fragile]{Atividade prática no laboratório}
  Escreva e execute:
  \begin{enumerate}
    \item Programa que lê \textbf{base} e \textbf{altura} e imprime a área do triângulo (formato com 2 casas)
    \item Programa que lê uma temperatura em Celsius e imprime em Fahrenheit: \(F = \tfrac{9}{5}C + 32\)
    \item Programa que lê o número de segundos e mostra em horas, minutos e segundos (use \texttt{//} e \texttt{\%})
    \item Programa que calcula a área de uma esfera de raio 5: \(A=4\pi r^2\) \hfill {\footnotesize (Exercício 2.18)}
  \end{enumerate}
  \begin{exampleblock}{Dica para o item 3}
    \texttt{h = seg // 3600} \hspace{1em} \texttt{resto = seg \% 3600} \hspace{1em} \texttt{m = resto // 60}
  \end{exampleblock}
\end{frame}

\exercicios{2}{%
  \begin{itemize}\small
    \item \textbf{2.9, 2.10} (área, segundos após a meia-noite)
    \item \textbf{2.17, 2.18} (hipotenusa, área da esfera)
  \end{itemize}
  \vspace{0.5ex}
  \footnotesize Capítulo 3: \href{\siteurl exercicios/cap3/}{\texttt{exercicios/cap3/}} -- \textbf{3.1, 3.2}.}

\resumo{
  \item Um programa estruturado combina sequência, seleção e iteração
  \item \texttt{print} formata a saída; \texttt{input} lê \emph{texto}, sempre
  \item Converta com \texttt{int()} e \texttt{float()} antes de calcular
  \item Parênteses tornam as expressões claras e seguras
  \item Próxima aula: estruturas de decisão (\texttt{if}, \texttt{elif}, \texttt{else})
}

\end{document}
